\section{Results and Discussion}
Table~\ref{tab:main-results} reports the current held-out test
metrics. XGBoost achieves a PR-AUC of 0.35409 and ROC-AUC of 0.62383,
slightly above Logistic Regression at 0.34915 and 0.61926. XGBoost
also has lower Brier score and ECE in this run. These differences
are modest and do not establish that the model will generalize
equally well to other periods or operational settings.

\begin{table}[htbp]
\centering
\caption{Current test-set results}
\label{tab:main-results}
\begin{tabular}{lrr}
\toprule
Metric & Logistic Regression & XGBoost \\
\midrule
PR-AUC $\uparrow$ & 0.34915 & 0.35409 \\
ROC-AUC $\uparrow$ & 0.61926 & 0.62383 \\
Brier score $\downarrow$ & 0.20013 & 0.19888 \\
ECE $\downarrow$ & 0.11950 & 0.11716 \\
Test precision at selected threshold & 0.50495 & N/A \\
Test recall at selected threshold & 0.00731 & N/A \\
\bottomrule
\end{tabular}
\end{table}

For Logistic Regression, the validation-selected threshold is
0.28551. On the test set it yields 102 true positives and 100 false
positives, with 13,844 false negatives and 39,945 true negatives.
Only 202 test observations are predicted positive. Although observed
precision is approximately 50.50\%, recall is approximately 0.73\%,
so this operating point identifies very few delayed flights.

For XGBoost, no validation threshold met the target precision of
60\%. Its threshold-dependent test metrics are therefore unavailable.
This is a failure to find a qualifying operating point under the
specified policy, not evidence of zero recall or zero precision.

The test-set positive rate is 25.83\%. Both models have moderate
ranking performance and calibration errors that leave room for
improvement. Threshold selection, calibration, and the cost of false
positives versus false negatives require further study.

Earlier reporting-buffer sensitivity experiments showed small
differences across buffer settings. Because the test set was evaluated
across those configurations, those comparisons are exploratory and
must not be used for further test-set tuning. The reporting buffer
is an assumption about historical outcome availability; its
operational validity has not been independently verified.
